\section{Expected Insights and Recommendations}

\subsection{High-Priority Signals}
The strongest expected signals are related to customer experience:
\begin{itemize}
  \item Low customer satisfaction.
  \item High number of complaints.
  \item High number of service calls.
\end{itemize}
These features are directly actionable because they can guide the customer
support and retention teams toward customers who may be experiencing service
problems.

\subsection{Moderate Signal}
Late payments are expected to have a moderate positive relationship with churn.
This feature should be interpreted together with satisfaction, complaint
history, and service calls before it is used for retention targeting.

\subsection{Weak Standalone Signals}
Average monthly GB usage, days since last interaction, and credit score are
expected to have weak standalone separation between churn and non-churn
customers. They may still be useful as filters or as supporting variables in a
multivariate model.

\subsection{Prediction Workflow Insight}
The Random Forest workflow in the group repository supports the EDA findings by
using interpretable churn-related features such as contract, satisfaction,
complaints, service calls, late payments, monthly charges, service count, and
tech support. Because the churn class is imbalanced, the model is more suitable
for ranking retention priority than for making a final yes/no business decision
without human review.

\subsection{Retention Recommendations}
\begin{enumerate}
  \item Prioritize customers with low satisfaction scores for proactive support.
  \item Create a complaint escalation view for customers with multiple
  complaints.
  \item Monitor repeated service calls as a sign of unresolved technical issues.
  \item Combine late-payment signals with experience-related signals before
  launching retention campaigns.
  \item Avoid drawing strong conclusions from usage volume or credit score
  alone.
\end{enumerate}

\subsection{Storytelling Flow for the Final Report}
The final report narrative should move from context to action:
\begin{enumerate}
  \item Establish why churn matters for a subscription-based telecom company.
  \item State the baseline churn rate and total customer population.
  \item Separate segment size from within-segment churn risk.
  \item Highlight above-baseline customer groups using EDA evidence.
  \item Connect the EDA findings with model evaluation and feature importance.
  \item End with retention actions that are tied to customer experience
  variables.
\end{enumerate}

This narrative avoids treating the model as the only result. The model supports
prioritization, while the visualization explains which segments are actionable
and why they deserve attention.
